\chapter{C- en C-achtige compilers}

Kort antwoord: CC65 en KickC werken beide op de MEGA65.

Zowel CC65 als KickC werken op de MEGA65, maar beide hebben standaard
alleen een C64-geheugenmodel en gebruiken uitsluitend 6502-opcodes.
Het zou geweldig zijn als iemand een C65-geheugenconfiguratie voor
CC65 zou ontwikkelen; dat zou niet al te moeilijk moeten zijn.

CC65 ondersteunt overlays, wat krachtig kan worden ingezet met het
extra geheugen van de MEGA65 om programma's groter dan 64 KB mogelijk
te maken. Daarvoor is echter een geschikte loader nodig, die nog niet
bestaat.

Op dezelfde manier zou het aanpassen van de codegenerator van CC65 om
45GS02-mogelijkheden te gebruiken niet bijzonder moeilijk zijn en zou
het helpen om de anders vreselijk trage en omvangrijke code die CC65
produceert te overwinnen. Het toevoegen van volledige ondersteuning
voor de 45GS02-CPU-mogelijkheden in CA65 (of nog beter: het laten
produceren van ACME-compatibele assembleeruitvoer door CC65) zou van
overweldigend voordeel zijn en niet bijzonder moeilijk te realiseren.
Dit zouden allemaal geweldige taken zijn om mee bezig te zijn terwijl
u wacht op uw MEGA65 DevKit!

Een voorbeeldsjabloon voor een C-programma dat met CC65 kan worden
gecompileerd en op de MEGA65 kan worden uitgevoerd, staat in het
repository \url{https://github.com/MEGA65/hello-world}. Dit repository
kan zelfs CC65 downloaden en compileren als u dat nog niet op uw
systeem hebt geïnstalleerd. Het repository werkt op Linux en Mac en op Windows via het Windows Subsystem for Linux (WSL).

Het LLVM-MOS-project (\url{https://llvm-mos.org}) heeft een
MEGA65-doelplatform dat is getest met C99, C++11 en Rust. De LLVM
MOS-compiler heeft ondersteuning voor 45GS02-instructies
(\texttt{-mcpu=mos45gs02}) en wordt geleverd met een assembler en
disassembler. Voorgecompileerde binaire bestanden van de SDK zijn
beschikbaar voor Linux, Windows en macOS.

\section{MEGA65 libc}

Er wordt een C-bibliotheek ontwikkeld voor de MEGA65, die al een
aantal nuttige functies bevat. Deze bibliotheek is beschikbaar via
\url{http://github.com/mega65/mega65-libc}. De procedures, functies
en definities die het biedt, zijn gedocumenteerd in een apart
hoofdstuk.

De MEGA65 libc is momenteel alleen beschikbaar voor LLVM en CC65,
hoewel we een KickC-port hiervan zeer op prijs zouden stellen.
